\section{Implementation and decision procedure}\label{sec:implementation}
\subsection{Input isolation and candidate work}
The research code separates model-visible observations from evaluator-only incident assignments and expert ranks. A receiver ledger uses availability rather than event timestamps to determine the visible set at $T=100$. Initial components are built from same-source, process-tree and eligible lineage references. Candidate relationship tokens bundle process, authentication and network observations by asserted endpoint pair. The verifier acts only on additional cross-fragment unions; it does not retroactively certify every pair in the initial components. This asymmetry is material to interpreting zero observed false grouped pairs.

The frozen holdout summary records 2,685 candidate pairs in each arm, 683 baseline accepted unions and 711 candidate accepted unions. The candidate's 28 additional accepts are precisely the fallback unions. The final partition contains 1,333 baseline groups and 1,305 candidate groups. The numerical difference in groups equals 28, but that equality does not prove a general one-union/one-group law: a valid token can connect alerts already in one component or fail a component guard in other populations. The code and audit distinguish token checks, accepted unions, and all within-group alert pairs.

\subsection{Union logic and complexity}
Let $\operatorname{find}(a)$ and $\operatorname{find}(b)$ identify current components. For each token in sorted order, evidence checks first decide whether the token is admissible; a union occurs only when the roots differ and the combined component satisfies the shared host/account and $100$-alert size guards. The receiver-visible original claims and endpoint fields are immutable inputs to each token check; an earlier union cannot manufacture an independent witness for a later token. Union--find gives amortized $O(\alpha(n))$ component operations after candidate generation, but total pipeline complexity depends on the number of token bundles, claims per bundle, and sorting. For $m$ bounded tokens and $n$ alerts, sorting costs $O(m\log m)$ and component operations $O(m\alpha(n))$; claim validation adds $O(\sum_r |B_r|)$ under a direct bundle scan. This is a structural bound on the implementation, not a measured throughput guarantee.

The candidate fallback has a strict branch condition. It is attempted only after the base verifier returns \texttt{missing\_family}; the process witness must be an original endpoint observation with a distinct origin root, and original identity and network claims must be timely and independently rooted. Any contradictory claim, mismatched causal field or component guard still blocks the union. Let $V_3(r)$ denote the baseline three-family predicate and $F_{2+e}(r)$ the guarded two-family-plus-endpoint predicate. The additional union decision is
\begin{equation}\label{eq:union}
U(r)=\bigl(V_3(r)\lor [\operatorname{reason}_3(r)=\mathrm{missing\_family}\land F_{2+e}(r)]\bigr)
\land C(r,\mathcal G_T),
\end{equation}
where $C$ checks the \emph{current} partition's host, account and size constraints. No oracle incident identifier enters $V_3$, $F_{2+e}$ or $C$. Equation~\ref{eq:union} clarifies why a rejected token stays a rejection in the truth audit even when it refers to a genuine relationship.

\subsection{Ranking coupling}
The group scorer in Eq.~\ref{eq:score} is identical in form for both holdout arms. The endpoint fallback can alter a component's members and accepted-token count, so it can change its score and the ranks of other groups. Group recoverability is therefore evaluated independently of retrieval at rank 20. The baseline and candidate share eligible alerts, seed, evaluator and ranking coefficients; the treatment is the narrow fallback rule with its consequent partition and token-bonus effects. This is a system-level comparison, not an isolated ranking ablation.

The archived runtime values, 0.696 versus 0.656 seconds, are single execution measurements. They do not establish a latency improvement: order, cache state, hardware load and repeated-run variability were not controlled for a throughput study. No measured P50/P95, analyst time or provider token cost exists for the frozen holdout.
